\documentclass[12pt]{article}
\usepackage[margin=1in]{geometry}
\usepackage[utf8]{inputenc}
\usepackage{latexsym, amsfonts, enumitem, graphicx, environ,enumitem, fancyhdr, color, amssymb, amsthm, amsmath, mathtools, tikz, nicematrix}
\usetikzlibrary{patterns}
\usetikzlibrary{decorations.pathreplacing}
\pagestyle{fancy}
\setlength{\headheight}{65pt}
\newenvironment{problem}[2][Problem]{\begin{trivlist}
    \item[\hskip \labelsep {\bfseries #1}\hskip \labelsep {\bfseries #2.}]}{\end{trivlist}}
\DeclareMathOperator{\Ima}{Im}

\usepackage{hyperref}
\usepackage[capitalize]{cleveref}

\newtheorem{thm}{Theorem}[section]
\newtheorem{lem}[thm]{Lemma}
\newtheorem{cl}[thm]{Claim}
\newtheorem{prop}[thm]{Proposition}
\newtheorem{cor}[thm]{Corollary}
\newtheorem{conj}[thm]{Conjecture}
\newtheorem{obstruction}[thm]{Obstruction}

\usepackage{mdframed}

\theoremstyle{definition}
\newtheorem{definition}[thm]{Definition}
\newtheorem{remark}[thm]{Remark}
\newtheorem{question}{Question}
\newtheorem{example}[thm]{Example}
\newtheorem{exercise}[thm]{Exercise}
\newtheorem{properties}[thm]{Properties}

\usepackage{tcolorbox}
\tcbuselibrary{theorems,skins,breakable}

\tcolorboxenvironment{example}{
enhanced jigsaw,colframe=cyan,interior hidden,
breakable,%before skip=10pt,after skip=10pt
}

\tcolorboxenvironment{thm}{
enhanced jigsaw,colframe=red,interior hidden,
breakable,%before skip=10pt,after skip=10pt
}

\tcolorboxenvironment{prop}{
enhanced jigsaw,colframe=red,interior hidden,
breakable,%before skip=10pt,after skip=10pt
}

\tcolorboxenvironment{properties}{
enhanced jigsaw,colframe=red,interior hidden,
breakable,%before skip=10pt,after skip=10pt
}

\tcolorboxenvironment{lem}{
enhanced jigsaw,colframe=red,interior hidden,
breakable,%before skip=10pt,after skip=10pt
}
\tcolorboxenvironment{cor}{
enhanced jigsaw,colframe=red,interior hidden,
breakable,%before skip=10pt,after skip=10pt
}
\tcolorboxenvironment{obstruction}{
enhanced jigsaw,colframe=red,interior hidden,
breakable,%before skip=10pt,after skip=10pt
}

\tcolorboxenvironment{proof}{
enhanced jigsaw, frame hidden,%interior hidden,
breakable,%before skip=10pt,after skip=10pt
}

\lhead{Avi Chetlin \\ Assignment 01 \\ 31 August 2026}
\rhead{Dr. Stanislaw Szarek \\ MATH 423 Measure Theory \\ Fall 2026}

\begin{document}
For a set \( A \subset [a,b] \), define the outer Riemann measure
\begin{equation}
\label{eq:1}
m^{*}(A) = \inf_{\cup J_i \supset A} \sum\limits_i l(J_i),
\end{equation}
where \( \left\{ J_i \right\} \) is a finite family of closed intervals contained in \( [a,b] \) and \( l(J) \) is the length of the interval \( J \).
One similarly defines the inner measure as
\begin{equation}
\label{eq:2}
m_{*}(A) = \sup_{\substack{\cup J_i \subset A \\ J_i \, \text{disjoint}}} \sum\limits_i l(J_i).
\end{equation}
If \( m^{*}(A) = m_{*}(A) \), then \( A \) is Riemann measureable and write \( A \in \mathcal{R}([a,b]) \).

\begin{problem}{1}
  Show that the value of \( m^{*} \) remains unchanged if
  \begin{enumerate}[label=(\roman*)]
    \item we require that the intervals \( (J_i) \) in the definition \cref{eq:1} are disjoint
          \begin{proof}
            Suppose that, with \( m^{*} \) defined as in \cref{eq:1}, a set \( A \) has outer measure \( m^{*}(A) = c \in [0, \infty] \).
            Let \( \hat{m}^{*}\) be the measure defined to require disjoint intervals.\\

            First of all, since any finite collection of disjoint closed intervals covering \( A \) is, in particular, a finite collection of (not necessarily disjoint) closed intervals covering \( A \), \( \hat{m}^{*} \) takes the infimum over a subset of that of \( m^{*} \), hence by elementary properties of the infimum,
            \begin{equation}
            \label{eq:6}
            m^{*} \leq \hat{m}^{*}.
            \end{equation}

            Conversely, let \( \left\{ K_i \right\} \) be an arbitrary finite collection of closed intervals covering \( A \).
            We apply the following algorithm:
            \begin{enumerate}[label=(\arabic*)]
              \item Consider \( K_1 \) and \( K_2 \); if \( K_1 \cap K_2 \neq \emptyset \), let \( \tilde{K}_1 = K_1 \cup K_{2} = [\min\{\alpha_1, \alpha_2\}, \max\{\beta_1, \beta_2\}] \).
              \item Consider \( \tilde{K}_1\) and \(K_3 \), proceed identically to step 1. Continue when \( \tilde{K}_1 \) and \( K_n \) are disjoint for some \( n \).
              \item Starting with \( K_n \), repeat steps 1 and 2 until you are left with a collection of disjoint closed intervals which still cover \( A \) (in particular the unions \( \bigcup  \tilde{K}_{i}  = \bigcup  K_i  \)).
            \end{enumerate}
            This algorithm necessarily halts after finite steps given the finite family of intervals.

            Now, since for any interval \( I = [\alpha, \beta]\), \( l(I) = \beta - \alpha \), it is clear that for two intersecting intervals \( I_1 \cap I_2 \neq \emptyset\),
            \begin{equation}
            l(I_1 \cup I_2) = l(\min\{\alpha_1, \alpha_2\}, \max\{\beta_1,\beta_2\}) = \max\{\beta_1,\beta_2\} - \min\{\alpha_1, \alpha_2\}.
            \end{equation}

            Without loss of generality, assume that \( \beta_1 > \beta_2 \), \( \alpha_{2} > \alpha_{1} \), so that \( l(I_1 \cup I_2) = \beta_1 - \alpha_{1} \). However,
            \begin{equation}
            l(I_1) + l(I_2) = \beta_1 - \alpha_1 + \beta_2 - \alpha_2 = (\beta_1 - \alpha_1) + (\beta_2 - \alpha_2).
            \end{equation}

            By definition of an interval, we must have that \( \beta_2 > \alpha_2 \), (even if the min/max were different, we always have \( \alpha_i < \beta_j \) since the intervals intersect), meaning that \( \beta_2 - \alpha_2 > 0 \), and hence
            \begin{equation}
            \label{eq:7}
            l(I_1 \cup I_2) \leq l(I_1) + l(I_2).
            \end{equation}

            Applying the algorithm, summing over the finite collection, and harnessing \cref{eq:7}, then taking the infimum of both sides yields
            \begin{equation}
            \label{eq:8}
            \hat{m}^{*} \leq m^{*}.
            \end{equation}
            Then \cref{eq:6,eq:8} show that \( \hat{m}^{*} = m^{*} \).

          \end{proof}

    \item we require that the family \( (J_i) \) is minimal (i.e., no proper subfamily of \( (J_i) \) covers \( A \)).
          \begin{proof}
            Define \( \hat{m}^{*} \) as \( m^{*} \) with the additional requirement of minimal families.
            Similarly to part (i), it is clear that since \( \hat{m}^{*} \) takes an infimum over a subset, we immediately have
            \begin{equation}
            \label{eq:9}
            m^{*} \leq \hat{m}^{*}.
            \end{equation}

            For the reverse inequality, note that for any arbitary, finite closed family \( \left\{ K_i \right\} \) which covers \( A \), if the family is non-minimal we can simply replace the family with the proper subfamily which still covers \( A \), \( \left\{ \tilde{K}_i \right\} \).

            Since this is a proper subfamily, that is, \( \left\{ \tilde{K}_i \right\} \subset \left\{ K_i \right\} \), we clearly have that
            \begin{equation}
            \label{eq:10}
            \sum\limits_i l(\tilde{K}_i) \leq \sum\limits_i l(K_i),
          \end{equation}
          (since the LHS has fewer summands than the RHS and \( l(I) \) is always non-negative.)

          Taking the infimum of both sides shows that
          \begin{equation}
          \label{eq:11}
          \hat{m}^{*} \leq m^{*},
          \end{equation}
          hence \( \hat{m}^{*} = m^{*} \).
          \end{proof}
    \item we require that the intervals \( J_i \) are open (rather than closed).
          \begin{proof}
            On one hand, note that for any finite family of open intervals \( \left\{ (\alpha_i, \beta_i) \right\} \) covering \( A \), can be associated uniquely with a finite closed cover \( \left\{ [\alpha_i, \beta_i] \right\} \) by taking the closure of each interval.
            Now,
            \begin{equation}
            \label{eq:13}
            l((\alpha_i, \beta_i)) = l([\alpha_i, \beta_i]),
          \end{equation}
          by definition.
          Since \( m^{*}(A) \) is a lower bound, we have that
          \begin{equation}
          \label{eq:14}
          m^{*} \leq \sum\limits_i \left\{ [\alpha_i, \beta_i] \right\} = \sum\limits_i \left\{ (\alpha_i, \beta_i) \right\}.
          \end{equation}
          That is, \( m^{*} \) is a lower bound for the sums of lengths of finite open covers.
          However, since \( \hat{m}^{*} \) is, by definition, the greatest lower bound, we must have that
          \begin{equation}
          \label{eq:15}
          m^{*} \leq \hat{m}^{*}.
          \end{equation}

          Conversely, consider a finite closed cover \( \left\{ [\alpha_i, \beta_i] \right\} \) such that, for some \(\varepsilon > 0\)
          \begin{equation}
          \label{eq:16}
          \sum\limits_i [\alpha_i, \beta_i] < m^{*}(A) + \frac{\varepsilon}{2},
          \end{equation}
          and expand each interval in the cover to an open interval
          \begin{equation}
          \label{eq:17}
          \left( \alpha_i - \delta, \beta + \delta \right)
          \end{equation}
          for some \(\delta > 0\).

          Now, the total length of this open cover is
          \begin{equation}
          \label{eq:18}
          \sum\limits_i (\alpha_i - \delta, \beta_i + \delta) = \sum\limits_i [\alpha_i, \beta_i] + 2n\delta
          \end{equation}
          where \(n\) is the total number of intervals in the cover.

          If we set \( \delta = \frac{\varepsilon}{4n} \), then
          \begin{equation}
          \label{eq:19}
          \sum\limits_i (\alpha_i - \delta, \beta_i + \delta) = \sum\limits_i [\alpha_i, \beta_i] + \frac{\varepsilon}{2} < m^{*}(A) + \varepsilon.
          \end{equation}

          Since this is true for arbitrary \(\varepsilon > 0\), we have that
          \begin{equation}
          \label{eq:20}
          \sum\limits_i (\alpha_i - \delta, \beta_i + \delta) \leq m^{*}(A),
          \end{equation}
          and therefore
          \begin{equation}
          \label{eq:21}
          \hat{m}^{*} \leq m^{*}.
          \end{equation}
          With \cref{eq:21,eq:15}, our claim follows.
          \end{proof}
  \end{enumerate}
\end{problem}

\begin{problem}{2}
  \label{prob:2}
  Show that \( m^{*}(A) = U(\chi_A) \) (see the handout) where \(\chi_A\) is the indicator function of \( A \).
  \textit{Note}: This is relatively subtle;
  start by associating with a partition of \( [a,b] \) the family of intervals of that partition that intersect \( A \).
\end{problem}
\begin{proof}
  We proceed by showing a bijection between the set of all finite closed disjoint covers of \( A \) and the set of all partitions of \( [a,b] \).

  First of all, given an arbitrary partition \( P = \{x_0, \dots, x_n\} \) of \( [a,b] \), let \( \{I_j\} \) be the set of intervals \( [x_{j-1}, x_{j}] \) from \( P \) such that \( [x_{j-1}, x_{j}] \cap A \neq \emptyset \).

  By definition,
  \begin{equation}
  \label{eq:23}
  U(\chi_A, P) = \sum\limits_{i=1}^n M_i \Delta x_i,
  \end{equation}
  where \( M_i = \sup_{[x_{i-1}, x_i]} \chi_A \) and \( \Delta x_i = x_i - x_{i-1} \).
  Since \( M_i = 1 \) for all \( i \) such that \( [x_{i-1}, x_i] \cap A \neq \emptyset\), it is clear that, (with \( j \) as above),
  \begin{equation}
  \label{eq:24}
  U(\chi_A, P) = \sum\limits_{j} \Delta x_i = \sum\limits_{j} l([x_{j-1}, x_j]).
  \end{equation}

  On the other hand, given an arbitrary finite closed disjoint cover \( \{ I_j \} \), define a partition \( P \) in the following manner.
  Let \( x_0 = a \), \( x_n = b \).
  Then, working iteratively, let \( x_1 = \alpha_1 \), the lower bound of \( I_1 \), \( x_2 = \beta_2 \), the upper bound of \( I_1 \), etc.
  Since we chose a disjoint cover (which by problem 1 is equivalent) and since the cover is contained in \( [a,b] \), this uniquely determines a partition of \( [a,b] \).

  Now, since the intervals of \( P \) with \( M_i = 1 \) are precisely those from the cover \( \{I_j\} \), it is clear that
  \begin{equation}
  \label{eq:25}
  \sum\limits_j l(I_j) = U(\chi_A, P).
  \end{equation}

  Now that we have established the bijection, we can take the infimum of both sets.
  For the set of partitions, that is
  \begin{equation}
  \label{eq:26}
  \inf_P U(\chi_A, P) = U(\chi_{A}),
\end{equation}
and for the set of covers,
\begin{equation}
\label{eq:27}
\inf_{\{I_j\}} \sum\limits_i l(I_j) = m^{*}(A).
\end{equation}

Since we established the bijection, the two sets are the same, and thus
\begin{equation}
\label{eq:28}
U(\chi_A) = m^{*}(A).
\end{equation}
\end{proof}

\begin{problem}{3}
  One can show that \( m_{*}(A) = L(\chi_A) \).
  Explain why this implies that \( A \) is Riemann measureable if and only if \(\chi_A \in \mathcal{R}[a,b]\), and that then \( m(A) = \int\limits_a^b \chi_A \).
\end{problem}
\begin{proof}
  Since in Problem 2 we showed that \( m^{*}(A) = U(\chi_A) \) and the definition of Riemann measureability is \( m^{*}(A) = U(\chi_{A}) = L(\chi_A) = m_{*}(A)\), we have precisely the definition of Riemann integrability for \( \chi_A \).

  Moreover, by definition, since \( U(\chi_A) = L(\chi_A) \), the integral of \( \chi_A \) is
  \begin{equation}
  \int\limits_a^b \chi_A = U(\chi_A) = L(\chi_A) = m_{*}(A) = m^{*}(A) = m_A,
  \end{equation}
  as desired.
\end{proof}

\begin{problem}{4}
  Show that if \( m^{*}(A) = 0 \), then \( A \) is Riemann measureable.
\end{problem}
\begin{proof}
  From problems 3 and 4, we know that \( m^{*}(A) = U(\chi_A) \) and \( m_{*}(A) = L(\chi_{A}) \).
  Recall from real analysis that for any function \( f \), \( U(f) \geq L(f) \).
  Also note that by definition, \( m_{*} \geq 0 \).
  We hence have that \( 0 = m^{*}(A) \geq m_{*}(A) \geq 0 \) which forces \( m_{*}(A) \).
  Thus, \( m^{*}(A) = m_{*}(A) = 0 \), so \( A \) is Riemann measureable.
\end{proof}

\begin{problem}{5}
  Suppose \(\mu: \mathcal{P}(\mathbb{R}) \rightarrow [0, \infty]\) is a finitely additive measure verifying
  \begin{enumerate}
    \item \(\mu(J) = l(J)\) for every interval \( J \subset \mathbb{R} \) and which is
    \item translation invariant.
  \end{enumerate}
  What is \(\mu(N)\), where \( N \subset [0,1) \) is the non-measureable set defined on p. 20 of the textbook?
\end{problem}
\begin{proof}
  Recall that \( N \) was defined as follows:
  Consider the set \( [0, 1) \), and an arbitrary element \( x \in [0,1) \).
  Define an equivalence relation \( \sim \) such that \( x \sim y  \) iff \( x - y \in \mathbb{Q} \).
  Each element is thus associated with its equivalence class \( [x] \), i.e. \(x \rightarrow (x + \mathbb{Q}) \cap [0,1) \).
  Then \( N \) is the set containing exactly one representative from each equivalence class.\\

  First of all, note that \( \mu([0,1)) = 1 \).
  Now, for each \( i \in \mathbb{N} \), define \( r_{i} \in \mathbb{Q} \cap [0,1) \) by choosing a unique element.
  Then write
  \begin{equation}
  \label{eq:29}
  N_{r_i} = \left\{ x + r_i \mid x \in N \cap [0,1-r_i) \right\} \cup \left\{ x + r_{i} - 1 \mid x \in N \cap [1-r_{i}, 1) \right\},
  \end{equation}
  as on p. 20 in the book.
  Define \( N^1 = N \), \( N^2 = N \cup N_{r_1} \), \( N^3 = N^2 \cup N_{r_{2}}\), and so on.
  Now, \( \mu(N^2) = \mu(N) + \mu(N_{r_1}) \) which by translation invariance is \(2\mu(N)\).
  Following the same line of argument inductively, we have that
  \begin{equation}
  \mu(N^n) = n\mu(N).
  \end{equation}

  Since \( N^{n} \subset [0,1) \) for all \( n \), we must have that \( \mu(N^n) \leq \mu([0,1)) = 1\).
  That is,
  \begin{equation}
  \label{eq:30}
  n\mu(N) \leq 1.
  \end{equation}
  For \cref{eq:30} to hold for all \( n \in \mathbb{N} \), we must have that \(\mu(N) = 0\).
  Thus,
  \begin{equation}
  \label{eq:31}
  \mu(N) = 0.
  \end{equation}
\end{proof}
\end{document}
